\documentclass[12pt,a4paper]{article}

\usepackage{fontspec}
\usepackage{polyglossia}
\setdefaultlanguage{greek}
\setotherlanguage{english}

\setmainfont{Times New Roman}
\setsansfont{Arial}
\setmonofont{Courier New}

\usepackage{geometry}
\usepackage{graphicx}
\usepackage{float}
\usepackage{caption}
\usepackage{xcolor}
\usepackage{listings}
\usepackage{amsmath}
\usepackage{hyperref}
\usepackage{fancyhdr}

\geometry{left=2.5cm,right=2.5cm,top=2.5cm,bottom=2.5cm}

\definecolor{codegreen}{rgb}{0,0.6,0}
\definecolor{codegray}{rgb}{0.5,0.5,0.5}
\definecolor{backcolour}{rgb}{0.95,0.95,0.92}

\lstset{
    backgroundcolor=\color{backcolour},
    commentstyle=\color{codegreen},
    keywordstyle=\color{blue},
    basicstyle=\ttfamily\small,
    breaklines=true,
    numbers=left,
    numberstyle=\tiny\color{codegray},
    language=Python,
    frame=single
}

\pagestyle{fancy}
\fancyhf{}
\rhead{AEM: 6931}
\lhead{Data Analysis 2025-26}
\cfoot{\thepage}

\title{
    \textbf{Ανάλυση Δεδομένων 2025-26}\\[0.3cm]
    \large Ερώτημα 1: Δειγματοληψία και Καθαρισμός Δεδομένων\\[0.2cm]
    \normalsize AEM: 6931
}
\author{}
\date{Μάρτιος 2026}

\begin{document}
\maketitle
\tableofcontents
\newpage

\section{Εισαγωγή}

\subsection{Σκοπός της Εργασίας}

Η παρούσα εργασία αφορά την ανάλυση δεδομένων ποιότητας κρασιού. Οι στόχοι είναι:

\begin{enumerate}
    \item Επιλογή τυχαίου δείγματος 80\% με seed τον ΑΕΜ
    \item Έλεγχος ποιότητας δεδομένων
    \item Καθαρισμός δεδομένων
\end{enumerate}

\subsection{Τι είναι ο Καθαρισμός Δεδομένων;}

\textbf{Ορισμός:} Ο καθαρισμός δεδομένων είναι η διαδικασία εντοπισμού και διόρθωσης λανθασμένων, ελλιπών, διπλότυπων ή ασυνεπών εγγραφών. Είναι κρίσιμο βήμα πριν από οποιαδήποτε ανάλυση.

Συνήθη προβλήματα:
\begin{itemize}
    \item \textbf{Τυπογραφικά λάθη:} π.χ. ``whte'' αντί ``white''
    \item \textbf{Ελλείπουσες τιμές:} Κενά πεδία
    \item \textbf{Μη έγκυρες τιμές:} π.χ. pH = 999
    \item \textbf{Διπλοεγγραφές:} Επαναλαμβανόμενες γραμμές
\end{itemize}

\section{Περιγραφή Δεδομένων}

\subsection{Μεταβλητές}

Το dataset περιέχει τις εξής μεταβλητές:

\begin{enumerate}
    \item \textbf{Fixed acidity} - Σταθερή οξύτητα (g/L)
    \item \textbf{Volatile acidity} - Πτητική οξύτητα (g/L)
    \item \textbf{Citric acid} - Κιτρικό οξύ (g/L)
    \item \textbf{Residual sugar} - Υπολειμματικά σάκχαρα (g/L)
    \item \textbf{Chlorides} - Χλωριούχα άλατα (g/L)
    \item \textbf{Free SO2} - Ελεύθερο διοξείδιο του θείου (mg/L)
    \item \textbf{Total SO2} - Συνολικό διοξείδιο του θείου (mg/L)
    \item \textbf{Density} - Πυκνότητα (g/cm3)
    \item \textbf{pH} - Οξύτητα (κλίμακα 2-14)
    \item \textbf{Sulphates} - Θειικά άλατα (g/L)
    \item \textbf{Alcohol} - Περιεκτικότητα σε αλκοόλ (\%)
    \item \textbf{Quality} - Βαθμολογία ποιότητας (0-10)
    \item \textbf{Wine\_type} - Τύπος κρασιού (red ή white)
\end{enumerate}

Αρχικό dataset: \textbf{8.385} γραμμές $\times$ \textbf{15} στήλες

\section{Δειγματοληψία}

\textbf{Seed = AEM = 6931} - Χρησιμοποιώντας τον ίδιο seed, παίρνουμε πάντα τα ίδια αποτελέσματα για αναπαραγωγιμότητα.

\begin{lstlisting}
import pandas as pd

df_original = pd.read_csv('Data Analysis_2026.csv')
AEM = 6931
df = df_original.sample(frac=0.80, random_state=AEM)
df = df.reset_index(drop=True)
\end{lstlisting}

Αποτέλεσμα: 8.385 $\rightarrow$ \textbf{6.708} γραμμές (80\%)

\section{Εντοπισμός Προβλημάτων}

\subsection{Τυπογραφικά Λάθη (wine\_type)}

Εντοπίστηκαν τα εξής λανθασμένα values στη στήλη wine\_type:

\begin{itemize}
    \item ``Redd '' (75 εγγραφές) $\rightarrow$ red
    \item ``whit'' (90 εγγραφές) $\rightarrow$ white
    \item ``r3d'' (64 εγγραφές) $\rightarrow$ red
    \item ``WHITTE'' (77 εγγραφές) $\rightarrow$ white
    \item ``whate'' (39 εγγραφές) $\rightarrow$ white
    \item 999, 9999, -999, ? (215 εγγραφές) $\rightarrow$ NaN
\end{itemize}

\textbf{Σύνολο τυπογραφικών λαθών:} 560

\subsection{Μη Αριθμητικές Τιμές}

Βρέθηκαν μη αριθμητικές τιμές (π.χ. ``abc'', ``?'', ``N/A'') στις αριθμητικές στήλες:

\begin{itemize}
    \item fixed acidity: 87 μη αριθμητικές τιμές
    \item volatile acidity: 45
    \item citric acid: 61
    \item pH: 151
    \item alcohol: 105
    \item κ.ά.
\end{itemize}

\textbf{Σύνολο μη αριθμητικών τιμών:} 845

\subsection{Αδύνατες Τιμές}

Τιμές όπως 999, -999, 9999 χρησιμοποιούνται συχνά ως κωδικοί για ελλείπουσες τιμές. Επίσης, κάποιες τιμές είναι εκτός λογικών ορίων (π.χ. pH = 999 ενώ το pH κυμαίνεται μεταξύ 0-14).

Λογικά όρια:
\begin{itemize}
    \item Fixed acidity: 0-20 g/L
    \item Volatile acidity: 0-5 g/L
    \item pH: 2-5
    \item Alcohol: 5-20\%
    \item Quality: 0-10
    \item Density: 0.9-1.1 g/cm$^3$
\end{itemize}

\textbf{Τιμές εκτός ορίων:} 2.037

\subsection{Διπλοεγγραφές}

\begin{lstlisting}
duplicates = df.duplicated()
print(f"Duplicates found: {duplicates.sum()}")
\end{lstlisting}

Εντοπίστηκαν \textbf{1.434} διπλοεγγραφές.

\subsection{Ελλείπουσες Τιμές (Missing Values)}

Μετά τη μετατροπή των μη αριθμητικών και αδύνατων τιμών σε NaN, οι ελλείπουσες τιμές ανά στήλη είναι:

\begin{itemize}
    \item fixed acidity: 574 (10.9\%)
    \item volatile acidity: 542 (10.3\%)
    \item density: 545 (10.3\%)
    \item pH: 644 (12.2\%)
    \item alcohol: 600 (11.4\%)
    \item quality: 40 (0.8\%)
\end{itemize}

\textbf{Σύνολο ελλειπουσών τιμών:} 3.092

\section{Αντιμετώπιση Προβλημάτων}

\subsection{Στρατηγική Αντιμετώπισης}

\begin{enumerate}
    \item \textbf{Τυπογραφικά λάθη wine\_type:} Διόρθωση στη σωστή τιμή ή NaN
    \item \textbf{Μη αριθμητικές τιμές:} Μετατροπή σε NaN
    \item \textbf{Αδύνατες τιμές:} Μετατροπή σε NaN
    \item \textbf{Διπλοεγγραφές:} Αφαίρεση
    \item \textbf{Missing values:} Συμπλήρωση με διάμεσο ανά τύπο κρασιού
\end{enumerate}

\subsection{Διόρθωση Τυπογραφικών Λαθών}

\begin{lstlisting}
corrections = {
    'Redd ': 'red', 'r3d': 'red',
    'WHITTE': 'white', 'whit': 'white', 'whate': 'white',
    '999': None, '9999': None, '-999': None, '?': None
}
for wrong, correct in corrections.items():
    df.loc[df['wine_type'] == wrong, 'wine_type'] = correct
\end{lstlisting}

\subsection{Γιατί Διάμεσος (Median) και όχι Μέσος Όρος;}

Η διάμεσος είναι η τιμή που χωρίζει τα δεδομένα στη μέση (50\% πάνω, 50\% κάτω). Είναι πιο ανθεκτική σε ακραίες τιμές (outliers) σε αντίθεση με τον μέσο όρο.

Παράδειγμα: Για τιμές [1, 2, 3, 4, 100]:
\begin{itemize}
    \item Μέσος όρος = 22 (επηρεάζεται από το 100)
    \item Διάμεσος = 3 (δεν επηρεάζεται)
\end{itemize}

\subsection{Γιατί ανά Τύπο Κρασιού;}

Τα κόκκινα και λευκά κρασιά έχουν διαφορετικά χημικά χαρακτηριστικά:

\begin{itemize}
    \item Κόκκινο - Fixed acidity (διάμεσος): 7.80
    \item Λευκό - Fixed acidity (διάμεσος): 6.80
    \item Κόκκινο - Total SO2 (διάμεσος): 40
    \item Λευκό - Total SO2 (διάμεσος): 133
\end{itemize}

\section{Αποτελέσματα}

\subsection{Συνοπτικός Πίνακας Προβλημάτων}

\begin{enumerate}
    \item Περιττές στήλες (Unnamed): 2 - Αφαιρέθηκαν
    \item Τυπογραφικά λάθη wine\_type: 560 - Διορθώθηκαν
    \item Μη αριθμητικές τιμές: 845 - Έγιναν NaN
    \item Αδύνατες τιμές: 2.037 - Έγιναν NaN
    \item Διπλοεγγραφές: 1.434 - Αφαιρέθηκαν
    \item Missing values (αρχικά): 3.092 - Συμπληρώθηκαν με διάμεσο
    \item Missing values (τελικά): 0
\end{enumerate}

\subsection{Εξέλιξη Μεγέθους Dataset}

\begin{enumerate}
    \item Αρχικό: 8.385 γραμμές × 15 στήλες
    \item Μετά sample (80\%): 6.708 γραμμές × 13 στήλες
    \item Μετά αφαίρεση duplicates: 5.274 γραμμές × 13 στήλες
    \item \textbf{Τελικό: 5.274 γραμμές × 13 στήλες}
\end{enumerate}

\subsection{Κατανομή Τελικού Dataset}

\begin{itemize}
    \item Λευκό κρασί: 3.924 (74.4\%)
    \item Κόκκινο κρασί: 1.350 (25.6\%)
    \item \textbf{Σύνολο: 5.274 (100\%)}
\end{itemize}

\section{Συμπεράσματα}

\begin{enumerate}
    \item Τα πραγματικά δεδομένα περιέχουν πολλαπλά προβλήματα ποιότητας
    \item Ο καθαρισμός δεδομένων είναι απαραίτητος πριν από κάθε ανάλυση
    \item Κάθε τύπος προβλήματος απαιτεί διαφορετική αντιμετώπιση
    \item Οι αποφάσεις καθαρισμού πρέπει να είναι τεκμηριωμένες
    \item Η χρήση seed εξασφαλίζει αναπαραγωγιμότητα
\end{enumerate}

Το καθαρισμένο dataset αποθηκεύτηκε στο αρχείο: \texttt{cleaned\_wine\_data\_AEM\_6931.csv}

\newpage
\appendix
\section{Πλήρης Κώδικας Python}

\begin{lstlisting}
import pandas as pd
import numpy as np

# Load data
df = pd.read_csv('Data Analysis_2026 1st Case_Data.csv')

# Remove unnamed columns
unnamed = [c for c in df.columns if 'Unnamed' in c]
df = df.drop(columns=unnamed)

# Sample 80%
AEM = 6931
df = df.sample(frac=0.80, random_state=AEM)
df = df.reset_index(drop=True)
print(f"After sampling: {len(df)} rows")

# Fix wine_type typos
fixes = {
    'Redd ':'red', 'r3d':'red', 
    'WHITTE':'white', 'whit':'white', 'whate':'white',
    '999':np.nan, '9999':np.nan, '-999':np.nan, '?':np.nan
}
for wrong, correct in fixes.items():
    df.loc[df['wine_type']==wrong, 'wine_type'] = correct

# Convert numeric columns
num_cols = ['fixed acidity', 'volatile acidity', 
            'citric acid', 'residual sugar', 'chlorides', 
            'free sulfur dioxide', 'total sulfur dioxide',
            'density', 'pH', 'sulphates', 'alcohol', 'quality']

for col in num_cols:
    df[col] = pd.to_numeric(df[col], errors='coerce')

# Replace impossible values with NaN
limits = {
    'fixed acidity': (0, 20),
    'volatile acidity': (0, 5),
    'citric acid': (0, 3),
    'residual sugar': (0, 100),
    'chlorides': (0, 1),
    'free sulfur dioxide': (0, 300),
    'total sulfur dioxide': (0, 500),
    'density': (0.9, 1.1),
    'pH': (2, 5),
    'sulphates': (0, 3),
    'alcohol': (5, 20),
    'quality': (0, 10)
}

for col, (min_val, max_val) in limits.items():
    mask = (df[col] < min_val) | (df[col] > max_val)
    df.loc[mask, col] = np.nan

# Remove duplicates
df = df.drop_duplicates().reset_index(drop=True)
print(f"After removing duplicates: {len(df)} rows")

# Fill missing values with median by wine type
df['wine_type'].fillna(df['wine_type'].mode()[0], inplace=True)

for col in num_cols:
    for wine in ['red', 'white']:
        mask = (df['wine_type'] == wine) & (df[col].isna())
        median = df.loc[df['wine_type'] == wine, col].median()
        df.loc[mask, col] = median

# Verify no missing values
print(f"Missing values: {df.isna().sum().sum()}")
print(f"Final shape: {df.shape}")

# Save cleaned data
df.to_csv('cleaned_wine_data_AEM_6931.csv', index=False)
print("Saved to cleaned_wine_data_AEM_6931.csv")
\end{lstlisting}

\end{document}
